% =============================================================================
% Odd Ratio Linear Fit - Comprehensive LaTeX Document for Papers
% =============================================================================
% This document provides a complete description of the odd ratio linear fit
% method, suitable for inclusion in the Methods section of scientific papers.
%
% Reference: Artigau et al. (2022), AJ, 164, 84, Appendix A
% Implementation: https://github.com/eartigau/odd_ratio_fits
% 
% Authors: Artigau, Cook, Cadieux, Stefanov (Atanas K.), Doyon
% =============================================================================

\documentclass[twocolumn]{aastex631}

\usepackage{amsmath}
\usepackage{amssymb}
\usepackage{graphicx}
\usepackage{hyperref}

\begin{document}

\title{Robust Linear Fitting with the Odd Ratio Mixture Model}

\author{\'Etienne Artigau}
\affiliation{D\'epartement de Physique, Universit\'e de Montr\'eal, Montr\'eal, QC H3C 3J7, Canada}
\affiliation{Institut de Recherche sur les Exoplan\`etes (iREx)}

\author{Neil J. Cook}
\affiliation{D\'epartement de Physique, Universit\'e de Montr\'eal, Montr\'eal, QC H3C 3J7, Canada}
\affiliation{Institut de Recherche sur les Exoplan\`etes (iREx)}

\author{Charles Cadieux}
\affiliation{D\'epartement de Physique, Universit\'e de Montr\'eal, Montr\'eal, QC H3C 3J7, Canada}
\affiliation{Institut de Recherche sur les Exoplan\`etes (iREx)}

\author{Atanas K. Stefanov}
\affiliation{Instituto de Astrof\'isica de Canarias, E-38205 La Laguna, Tenerife, Spain}

\author{Ren\'e Doyon}
\affiliation{D\'epartement de Physique, Universit\'e de Montr\'eal, Montr\'eal, QC H3C 3J7, Canada}
\affiliation{Institut de Recherche sur les Exoplan\`etes (iREx)}


\begin{abstract}
We present a robust linear fitting algorithm based on a Gaussian mixture model
approach for outlier rejection. The method extends the odd ratio weighted mean
introduced in Appendix A of \citet{Artigau2022} to linear and polynomial
regression. Unlike traditional sigma-clipping, this approach uses a probabilistic
model to iteratively down-weight outliers while preserving proper statistical
uncertainties. The algorithm correctly handles heteroscedastic data, identifying
outliers based on their deviation in units of their own uncertainty rather than
absolute value. Monte Carlo simulations confirm that the returned uncertainties
are statistically valid. The method is implemented in the open-source Python
package \texttt{odd\_ratio\_linfit}.
\end{abstract}

\keywords{methods: statistical --- techniques: radial velocities --- methods: data analysis}


% =============================================================================
\section{Introduction}
\label{sec:intro}
% =============================================================================

Robust fitting in the presence of outliers is a common challenge in astronomical
data analysis. Traditional approaches such as sigma-clipping use hard thresholds
to reject data points that deviate significantly from the model. However, this
approach has several drawbacks: (1) the threshold is arbitrary, (2) points near
the threshold may be incorrectly classified, and (3) the resulting error estimates
become biased due to the sample truncation.

We present an alternative approach based on a Gaussian mixture model, extending
the odd ratio weighted mean method introduced in Appendix A of \citet{Artigau2022}
to linear and polynomial regression. This method treats outlier identification as
a probabilistic inference problem, allowing points to be smoothly down-weighted
rather than hard-rejected.


% =============================================================================
\section{Method}
\label{sec:method}
% =============================================================================

\subsection{Mixture Model Formulation}
\label{sec:mixture}

Consider fitting a linear model $y = a + bx$ to data $\{x_i, y_i, \sigma_i\}_{i=1}^N$,
where $\sigma_i$ is the measurement uncertainty on $y_i$. We model each data point
as arising from a two-component mixture distribution:
\begin{equation}
    P(y_i | x_i, a, b) = (1-f_0) \cdot \mathcal{N}(y_i | a+bx_i, \sigma_i) + f_0 \cdot P_{\mathrm{out}}
    \label{eq:mixture}
\end{equation}
where $\mathcal{N}(\cdot|\mu,\sigma)$ denotes a Gaussian distribution with mean $\mu$
and standard deviation $\sigma$, $f_0$ is the prior probability that any given point
is an outlier, and $P_{\mathrm{out}}$ represents the outlier distribution, assumed to
be uniform and broad.

Given this mixture model, the posterior probability that point $i$ with residual
$r_i = y_i - (a + bx_i)$ is a valid (non-outlier) measurement is:
\begin{equation}
    p_{\mathrm{valid},i} = \frac{(1-f_0) \exp\left(-r_i^2/2\sigma_i^2\right)}
        {(1-f_0) \exp\left(-r_i^2/2\sigma_i^2\right) + f_0}
    \label{eq:pvalid}
\end{equation}

Equivalently, the probability that point $i$ is an outlier is:
\begin{equation}
    p_{\mathrm{out},i} = \frac{f_0}{(1-f_0) \exp\left(-r_i^2/2\sigma_i^2\right) + f_0}
    \label{eq:pout}
\end{equation}

The default value $f_0 = 2 \times 10^{-4}$ implies that roughly 1 in 5000 points is
expected to be an outlier \emph{a priori}. This value provides a good balance between
robustness and sensitivity for typical astronomical datasets. Points deviating by
$\gtrsim 3.5\sigma$ from the model are significantly down-weighted.


\subsection{Iterative Fitting Algorithm}
\label{sec:algorithm}

The fitting proceeds via an iterative expectation-maximization-like procedure:

\begin{enumerate}
    \item \textbf{Initialization}: Perform standard weighted least squares with
          weights $w_i = 1/\sigma_i^2$ to obtain initial estimates $\hat{a}$, $\hat{b}$.
    
    \item \textbf{E-step}: Compute $p_{\mathrm{valid},i}$ for each point using
          Equation~\ref{eq:pvalid} with current parameter estimates.
    
    \item \textbf{M-step}: Update the model parameters using weighted least squares
          with weights $w_i = p_{\mathrm{valid},i}/\sigma_i^2$:
          \begin{align}
              \hat{b} &= \frac{\sum_i w_i (x_i-\bar{x})(y_i-\bar{y})}{\sum_i w_i (x_i-\bar{x})^2} \\
              \hat{a} &= \bar{y} - \hat{b}\bar{x}
          \end{align}
          where $\bar{x} = \sum_i w_i x_i / \sum_i w_i$ and similarly for $\bar{y}$.
    
    \item \textbf{Convergence check}: If the relative change in parameters is below
          a threshold (default: $10^{-2}$ times the parameter uncertainty), stop;
          otherwise return to step 2.
\end{enumerate}

The algorithm typically converges in 3--5 iterations.


\subsection{Uncertainty Estimation}
\label{sec:errors}

Parameter uncertainties are derived from the weighted covariance matrix:
\begin{align}
    \sigma_a^2 &= \frac{\sum_i w_i x_i^2}{\Delta} \\
    \sigma_b^2 &= \frac{\sum_i w_i}{\Delta}
    \label{eq:errors}
\end{align}
where $\Delta = (\sum_i w_i)(\sum_i w_i x_i^2) - (\sum_i w_i x_i)^2$, and the weights
$w_i = p_{\mathrm{valid},i}/\sigma_i^2$ incorporate the probability that each point
is valid.

This formulation ensures that the uncertainties properly account for the effective
reduction in sample size due to down-weighting of outliers.


% =============================================================================
\section{Key Properties}
\label{sec:properties}
% =============================================================================

\subsection{Proper Handling of Heteroscedastic Data}
\label{sec:heteroscedastic}

A crucial advantage of this method is its correct treatment of heteroscedastic
(varying) uncertainties. Outliers are identified based on their deviation in units
of their own uncertainty ($\sigma$), not absolute deviation.

Consider two points with the same absolute deviation from the fit:
\begin{itemize}
    \item Point A with small uncertainty $\sigma_A = 0.3$ and deviation $|r_A| = 3$:
          This is a $10\sigma$ outlier, heavily down-weighted.
    \item Point B with large uncertainty $\sigma_B = 2.0$ and deviation $|r_B| = 3$:
          This is only a $1.5\sigma$ deviation, consistent with normal scatter, retained
          with high weight.
\end{itemize}

This is the statistically correct behavior: a measurement with large uncertainty that
happens to deviate from the fit is not necessarily an outlier---it may simply reflect
its larger measurement error.


\subsection{Statistically Valid Uncertainties}
\label{sec:validation}

The uncertainties returned by the algorithm are statistically meaningful and properly
calibrated. This is verified through Monte Carlo simulations (Figure~\ref{fig:mc}).

For 1000 realizations of linear fitting with 10\% outliers:
\begin{itemize}
    \item The actual scatter in recovered parameters matches the mean reported
          uncertainties (ratio $\approx 1.0$).
    \item Normalized residuals $(p - p_{\mathrm{true}})/\sigma_p$ follow a standard
          normal distribution $\mathcal{N}(0,1)$.
\end{itemize}

This confirms that the reported errors correctly describe the statistical uncertainty.


\subsection{Advantages over Sigma-Clipping}
\label{sec:advantages}

The odd ratio method offers several advantages over traditional sigma-clipping:

\begin{enumerate}
    \item \textbf{No arbitrary thresholds}: The single parameter $f_0$ has a clear
          probabilistic interpretation.
    \item \textbf{Smooth weighting}: Points are partially down-weighted rather than
          completely rejected, avoiding discontinuities.
    \item \textbf{Stable convergence}: Binary inclusion/rejection can cause oscillations
          in iterative algorithms, where points flip in and out of the fit. Smooth
          weights adjust gradually at each iteration, ensuring stable convergence.
    \item \textbf{Proper error propagation}: Uncertainties account for the effective
          sample size reduction.
    \item \textbf{Heteroscedastic handling}: Different measurement uncertainties are
          naturally incorporated.
\end{enumerate}


% =============================================================================
\section{Implementation}
\label{sec:implementation}
% =============================================================================

The method is implemented in the open-source Python package \texttt{odd\_ratio\_fits},
available at \url{https://github.com/eartigau/odd_ratio_fits}. The package provides
three main functions:

\begin{itemize}
    \item \texttt{orf.mean}: Robust weighted mean
    \item \texttt{orf.linear}: Robust linear regression $y = a + bx$
    \item \texttt{orf.polyfit}: Robust polynomial regression
\end{itemize}

Basic usage:
\begin{verbatim}
import odd_ratio_fits as orf

a, a_err, b, b_err = orf.linear(x, y, yerr)
\end{verbatim}

The package depends only on NumPy and is compatible with Python 3.8+.


% =============================================================================
\section{Limitations}
\label{sec:limitations}
% =============================================================================

The method assumes:
\begin{itemize}
    \item Outliers are rare ($\lesssim 20$--30\% of data)
    \item Valid measurements follow Gaussian errors
    \item The linear (or polynomial) model is appropriate
\end{itemize}

If outliers dominate the dataset, the algorithm may fail to converge to the correct
solution. In such cases, more robust initialization or a higher $f_0$ value may help.


% =============================================================================
\section{Conclusion}
\label{sec:conclusion}
% =============================================================================

We have presented a robust linear fitting method based on a Gaussian mixture model
for outlier rejection. The method correctly handles heteroscedastic data, provides
statistically valid uncertainties, and offers several advantages over traditional
sigma-clipping approaches. The implementation is available as an open-source Python
package for use by the astronomical community.


\acknowledgments

This work was supported by the Institut de Recherche sur les Exoplan\`etes (iREx).


\begin{thebibliography}{99}
\bibitem[Artigau et al.(2022)]{Artigau2022}
Artigau, \'E., Cadieux, C., Cook, N.~J., et al.\ 2022, \aj, 164, 84
\end{thebibliography}

\end{document}
